\section{Conclusion and Future Work}
\label{sec:conclusion}

This paper introduced \dreamprice{}, to our knowledge the first Dreamer-style world model trained directly on retail scanner data. The system combines \dreamer{}'s three-phase training recipe with a \mamba{} selective state-space backbone, DRAMA-style decoupled posteriors, and a causally-constrained demand decoder. Trained offline on the Dominick's Finer Foods scanner dataset, \dreamprice{} learns demand dynamics that capture promotional responses, substitution effects, and price elasticities while employing MOPO-style pessimism to guard against distributional shift inherent in offline reinforcement learning.

The contributions span five dimensions. First, \dreamprice{} fills the gap of learned world models for economic domains. Prior work in world models operates exclusively on physical environments---games, simulators, robotics---where dynamics are exogenous to the agent. In retail, prices and demand are simultaneously determined; the transition dynamics are endogenous. \dreamprice{} is the first system to learn such dynamics from scanner data, enabling imagination-based planning and counterfactual demand estimation in a domain previously served only by hand-crafted simulators or model-free methods.

Second, the \mamba{} backbone replaces the GRU recurrence of standard RSSMs with structured state-space duality. During training, the model processes sequences in parallel via the SSD scan; during imagination, it switches to recurrent \texttt{step()} mode for $O(1)$-per-step rollouts. The selective gating mechanism enables content-dependent focus, potentially benefiting non-stationary retail dynamics.

Third, the DRAMA-style decoupled posterior---where the stochastic latent depends only on the current observation, not on the recurrent hidden state---breaks the sequential dependency that would otherwise prevent \mamba{}'s parallel scan during training. This design choice, following \citet{rajbhandari2024drama}, is essential for scalability.

Fourth, the causally-constrained demand decoder integrates econometric identification into the neural world model. Per-category price elasticities are pre-estimated via DML-PLIV \citep{chernozhukov2018dml} and frozen into the decoder. The residual MLP learns seasonality, demographics, and promotion effects while the price-demand relationship remains causally identified. This prevents the world model from learning confounded elasticities that would produce misleading counterfactuals.

Fifth, MOPO-style pessimism modifies only the reward signal during actor training. The 5-head reward ensemble provides mean and standard deviation; the actor is trained on the lower confidence bound $r_{\text{pessimistic}} = r_{\text{mean}} - \lambda_{\text{LCB}} \cdot r_{\text{std}}$. This minimal modification preserves the three-phase structure while addressing the distribution shift problem of offline reinforcement learning \citep{levine2020offline, yu2020mopo}.

\paragraph{Limitations.} Several limitations constrain the scope of the results. The evaluation is confined to a single product category (canned soup). Generalization to other categories---beer, soft drinks, cereals---requires additional training and validation; the entity-factored representation may support transfer learning across categories, but this remains an open question. The system is offline-only: no real-time deployment or online fine-tuning is evaluated. The Dominick's data span 1989--1997. Retail dynamics have evolved substantially since then: e-commerce, dynamic pricing algorithms, and supply chain digitization may alter the structure of demand and competition. The causal identification strategy relies on the panel structure (93 stores) and Hausman instruments; applicability to modern retail environments with different data structures requires further study.

\paragraph{Future work.} Several directions extend this work. Inference-time price optimization via CEM or gradient-based planning would enable the world model to produce optimal price sequences at deployment rather than relying solely on the trained actor. Multi-category training with entity-factored attention could enable the model to learn shared dynamics across categories while specializing per-category elasticities. Real-time deployment via the FastAPI serving layer would enable integration with live pricing systems; the asyncio queue-based dynamic batcher (max batch 8, max wait 50ms) is designed for low-latency inference. Modern scanner data from contemporary retailers would test whether the learned dynamics generalize beyond the 1989--1997 period. Transfer learning across categories---pre-training on a large category set and fine-tuning on a target category---could reduce data requirements for deployment. Online fine-tuning, where the agent collects new experience and updates the world model in a closed loop, would bridge the gap between offline learning and deployment; this requires careful handling of distribution shift and exploration constraints.

\paragraph{Reproducibility and open-source release.} The complete source code, trained models, data, and interactive demonstrations are publicly available to facilitate reproduction and extension:
\begin{itemize}[nosep,leftmargin=*]
  \item \textbf{Code}: \url{https://github.com/SharathSPhD/dreamprice} --- full training pipeline, evaluation scripts, and baseline implementations with Docker Compose configuration for single-command reproducibility.
  \item \textbf{Dataset}: \url{https://huggingface.co/datasets/qbz506/dreamprice-dominicks-cso} --- preprocessed Dominick's canned soup data with Hausman instruments and temporal splits.
  \item \textbf{Model}: \url{https://huggingface.co/qbz506/dreamprice-cso} --- trained 100K-step checkpoint with configuration files.
  \item \textbf{Demo}: \url{https://huggingface.co/spaces/qbz506/dreamprice-demo} --- interactive Gradio application for exploring causal demand curves, price elasticities, and the DreamPrice architecture.
  \item \textbf{Experiment tracking}: \url{https://wandb.ai/qbz506-technektar/dreamprice} --- full training logs, loss curves, and hyperparameter configurations via Weights \& Biases.
\end{itemize}

\dreamprice{} establishes that learned world models can operate in economic domains. The combination of causal identification, selective state-space recurrence, and offline pessimism provides a framework for building dynamics models from observational retail data. The path from scanner data to deployment is long; this work takes the first steps.
